\begin{figure*}[ht!]
 \centering
 \includegraphics[width=\textwidth]{figures/evaluation_overview.pdf}
 \vspace{-7mm}
 \caption{Overview of \eval~ evaluation framework: (1) Firstly, LVLMs generate captions from \data~ benchmark images. (2) Following this, LLMs are employed to \textit{extract} pivotal features that encapsulate from the generated descriptions. (3) Subsequently, these features are \textit{aligned} with a pre-defined list of ground-truth features using LLMs, facilitating the creation of two essential outputs: a dictionary of matched features and a more extensive dictionary encompassing broader conceptual matches. (4) Finally, we calculate two key metrics: \textit{faithfulness} and \textit{coverage}. These metrics measure the LVLMs' comprehension by evaluating how well the generated captions encapsulate the salient features of the images and the breadth of concepts they cover, respectively.}
 \vspace{-2mm}
 \label{fig:evaluation_overview}
\end{figure*}

\begin{table*}[t]
    \small
    \centering
    \renewcommand{\arraystretch}{1.1}
    \setlength{\tabcolsep}{0.6pt}
    \begin{tabular}{l c c c c c c c c c c c c}
    \toprule    
    \multirow{4}{*}{\textbf{Model}}&\multicolumn{2}{c}{\textbf{Object}}&\multicolumn{4}{c}{\textbf{Attribute}}&\multicolumn{4}{c}{\textbf{Relation}}&\multirow{4}{*}{\makecell{Average \\ Faithful. \\Score \\ (\%)}}&\multirow{4}{*}{\makecell{Average \\ Cover. \\Score \\ (\%)}}\\
    \cmidrule(lr){2-3} 
    \cmidrule(lr){4-7}
    \cmidrule(lr){8-11}
    & \multicolumn{2}{c}{\textit{Existence}} & \multicolumn{4}{c}{\textit{Color \& Counting}} & \multicolumn{2}{c}{\textit{Positional}} & \multicolumn{2}{c}{\textit{Comparative}} \\
    \cmidrule(lr){4-7}

& \multicolumn{2}{c}{} & \multicolumn{2}{c}{Object} & \multicolumn{2}{c}{People} & \multicolumn{2}{c}{} & \multicolumn{2}{c}{} \\
    \cmidrule(lr){2-3} 
    \cmidrule(lr){4-5}
     \cmidrule(lr){6-7}
     \cmidrule(lr){8-9}
    \cmidrule(lr){10-11}
&Faithful$_{\uparrow}$&Cover$_{\uparrow}$&Faithful$_{\uparrow}$&Cover$_{\uparrow}$&Faithful$_{\uparrow}$&Cover$_{\uparrow}$&Faithful$_{\uparrow}$&Cover$_{\uparrow}$&Faithful$_{\uparrow}$&Cover$_{\uparrow}$&\\
    \midrule
    InstructBLIP & 74.5 & 24.8 & 72.0 & 23.9 & 47.1 & 9.3 & 50.0 & 13.6 & 66.9 & 35.6 & 62.1 & 21.44 \\
    LLaVA-1.5 & 72.1  & 24.7 & 74.6 & 37.8 & 43.3 & 12.1 & 64.8 & 14.9 & 51.9 & \hlc{lightblue}{40.1~} & 61.34 & 25.92 \\ 
    MiniGPT-4 v2  & 65.0 & 25.4 & \hlc{lightyellow}{64.5~} & 17.9 & 38.9 & 11.6 & \hlc{lightyellow}{38.8~} & \hlc{lightblue}{33.1~} & 44.7 & \hlc{lightyellow}{11.2~} & \hlc{lightyellow}{50.38~} & 19.84 \\
    mPLUG-Owl2 & 71.5 & 24.8 & \hlc{lightblue}{79.9~} & 32.7 & 39.7 & 16.2 & 45.2 & 10.8 & \hlc{lightyellow}{41.6~} & 30.6 & 55.58 & 23.02 \\
    BLIVA & 77.7 & 21.9 & 73.3 & 24.3 & 37.6 & 11.6 & 39.5 & 9.7 & 68.0 & 29.9 & 59.22 & 19.48 \\
    CogVLM  & 71.2 & 35.5 & 75.3 & 24.3  & 43.7 & 22.4 & 51.9 & 10.5 & 49.0 & 35.9 & 58.22 & 25.72 \\
    InternLM-XComposer2 & 82.5 & 23.9  & 75.8 & 26.3 & 50.4 & 13.8 & 62.6 & 11.1 & 64.1 & 38.4 & 67.08 & 22.7 \\
    Qwen-VL-Chat & 70.6 & 28.4 & 75.1 & \hlc{lightblue}{38.6~} & 38.8 & 16.0 & 56.9 & 8.5 & 51.9 & 24.3 & 58.66 & 23.16 \\ 
    Emu2 & \hlc{lightblue}{94.2~} & \hlc{lightyellow}{14.1~} & 66.7 & \hlc{lightyellow}{10.4~} & \hlc{lightblue}{54.3~} & \hlc{lightyellow}{1.9~} & \hlc{lightblue}{72.2~} & \hlc{lightyellow}{1.8~} & \hlc{lightblue}{87.5~} & 12.3 & \hlc{lightblue}{74.98~} & \hlc{lightyellow}{8.1~} \\
    GPT-4V & \hlc{lightyellow}{61.6} & \hlc{lightblue}{38.8~} & 78.5 & 36.3 & \hlc{lightyellow}{34.7~} & \hlc{lightblue}{23.8~} & 46.7 & 12.6 & ~~51.6$^*$ & ~~28.5$^*$ &54.62 & \hlc{lightblue}{28.0~} \\
    \bottomrule
    \end{tabular}
    \vspace{-2mm}
    \caption{The overall evaluation results of object existence, attribute, and relation hallucination in \data~ using GPT-4 as the LLM Agent within \eval~. The highest is highlighted in \hlc{lightblue}{blue}, while the worst performance is highlighted in \hlc{lightyellow}{yellow}. Faithfulness and coverage scores are in percentage (\%). For images that contain people, GPT-4V refrains from generating comments, and we marked this score with an asterisk ($*$).}
    \vspace{-4mm}
    \label{tab:evaluation_results}
\end{table*}

\section{\eval~} 
\label{sec:evaluation_framework}

We propose a framework \eval~ that generalizes CHAIR, a metric that is widely adopted in existing studies~\cite{Zhai2023HallESwitchRA, wang2023llm}, by introducing semantic matching and incorporating both the \textit{faithfulness} and \textit{coverage} aspects into the evaluation. As shown in Figure~\ref{fig:evaluation_overview}, our evaluation process has two steps: \textit{feature extraction} and \textit{matching} (\Cref{sec:features_extract_match}) and \textit{scoring} (\Cref{sec:evaluation_metrics}).

\vspace{-2mm}
\subsection{Feature Extraction and Matching}
\label{sec:features_extract_match}

We start the process by generating an initial response, denoted as $R$, using a specific LVLM with the input pair $(I, p_G)$, where $I$ denotes the image and $p_G$ represents the prompt designed for LVLMs generation from \data~. Then, we leverage an LLM to analyze $R$ and extract key features. This is achieved through a series of prompts $p_E$, outlined in \Cref{apx:features_extraction_prompts}, which are designed to \textit{extract} features from object existence, attributes, and relations, respectively, resulting in a comprehensive list of extracted features from $R$, denoted as $F_R=\{f_{R_1}, f_{R_2}, ..., f_{R_m}\}$. Next, we utilize an LLM to \textit{align} the extracted features list $F_R$ with a pre-annotated ground-truth features list $F_G=\{f_{G_1}, f_{G_2}, ..., f_{G_m}\}$ from \data~. This alignment is facilitated through a set of carefully crafted prompts $p_M$, outlined in \Cref{apx:features_matching_prompts}, tailored to each feature subset, aiming to identify correlations and correspondences. Unlike previous evaluation metrics that rely on a fixed feature list and direct mapping, our approach eschews pre-processing and instead utilizes LLMs' language comprehension capabilities to semantically match extracted features with their ground-truth counterparts. This process yields two key outputs: \textbf{matched features} dictionary ($D_{M}$) and \textbf{broader conceptual matches} dictionary ($D_{B}$).\looseness=-1

$D_{M}$ contains features $f_{R_{i'_m}}$ from $f_R$ that semantically aligned with the features $f_{G_{i_m}}$ from $F_G$, ensuring \textit{precision}. For example, if we have the extracted ``(plaid, shirts)'' and the candidate ground-truth feature is ``(checkered, shirt),'' we can establish a match between these two because ``plaid'' and ``checkered'' are conceptually similar patterns often used interchangeably in the context of textiles. 

$D_{B}$ includes features $f_{R_{j'_n}}$ from $f_R$ that have broader conceptual meanings than the features $f_{G_{j_n}}$ from $F_G$, adding \textit{conceptual depth} to the evaluation. For instance, if we have the extracted ``(red, clothes)'' from an image, and the ground-truth annotation is ``(red, dress),'' we can still consider these features to match. This is because ``clothes'' is a broader category that encompasses ``dress.'' Therefore, despite the slight difference in specificity, the extracted features can be aligned with the ground-truth annotations based on their semantic relationship, where ``dress'' is a sub-type of ``clothes.''

\subsection{Evaluation Metrics}
\label{sec:evaluation_metrics}

We introduce two metrics to evaluate the hallucinations in two dimensions: \textit{faithfulness} and \textit{coverage} based on the original CHAIR metric. 

\paragraph{Faithfulness.} In the context of image captioning, faithfulness measures how closely captions match an image's content, emphasizing \textit{accuracy} in depicting visual elements and their attributes and relations without introducing hallucinations. It is calculated by comparing generated features against actual image features, considering both \textit{direct} ($D_{M}$) and \textit{broader} conceptual similarities ($D_{B}$):
\begin{equation}
    \small
    \text{Faithfulness}(R, F_G) = \frac{|D_{M} \cup set(D_{B})|}{|F_R|} \in [0,1].
\end{equation}

\paragraph{Coverage.} It measures the \textit{comprehensiveness} of the generated captions in capturing the key elements and attributes depicted in the image. It evaluates the proportion of ground-truth features that are successfully captured in the generated response, only through \textit{direct} matches ($D_{M}$):
\begin{equation}
    \small
    \text{Coverage}(R, F_G) = \frac{|set(D_{M})|}{|F_G|} \in [0,1].
\end{equation}